% ---------------------------------------------------------------------------
% The one-pass move: one forward and one backward pass through ESM3 per Monte
% Carlo move, instead of three forwards and two backwards.
%
% Proposal, not yet implemented or tested. Companion to zero_polymer_recap.tex
% (where the cost measurements come from), informedness.tex (pointing
% probability, dt/sqrt(T) scaling) and gradient_method.tex (gradient variants,
% cost flat in the number of relaxed sites).
%
% Build:  cd softmax/docs && pdflatex one_pass_move.tex && pdflatex one_pass_move.tex
% ---------------------------------------------------------------------------
\documentclass[11pt,a4paper]{article}

\usepackage[T1]{fontenc}
\usepackage[utf8]{inputenc}
\usepackage[margin=2.5cm]{geometry}
\usepackage{amsmath,amssymb}
\usepackage{booktabs}
\usepackage{array}
\usepackage[hidelinks]{hyperref}
\sloppy

\newcommand{\dt}{\Delta t}
\newcommand{\Mass}{M}
\newcommand{\Tsftm}{T_{\mathrm{sftm}}}
\newcommand{\Uam}{U_{\mathrm{am}}}
\newcommand{\file}[1]{\texttt{#1}}

\title{The one-pass move\\
       \large One forward and one backward pass per Monte Carlo move,\\
       with the gradient taken at the exact sequence and reused}
\author{\texttt{my\_extendedSampling} / \texttt{softmax}}
\date{2026--10--22}

\begin{document}
\maketitle

\begin{abstract}
\noindent
A move of the current rate sampler costs two forward+backward passes through
ESM3 and one extra forward pass, $\approx361$\,ms on an RTX 4500 Ada, and the
GPU is fully busy while it runs. This note proposes a move that costs one
forward and one backward pass, $\approx150$\,ms, an expected speed-up of
$\approx2.4\times$ at equal acceptance. Two changes make it possible. The
gradient is taken with respect to the one-hot probability matrix at the exact
sequence, so the forward pass that computes the exact energy also yields the
gradient. And the gradient is computed for all sites at once, which costs the
same as for one site, so the gradient at the candidate becomes the gradient
of the next move whether the candidate is accepted or not. The kernel stays an
exact Metropolis--Hastings kernel for $e^{-U/T}$, including when the backward
pass is not bit-reproducible. The gradient changes scale but not its ranking of
candidates, so $\dt$ must be rescaled by a computable factor,
$c=p_{\mathrm{off}}/\Tsftm\approx4.5\times10^{-4}$. Nothing here has been run
yet; Section~\ref{sec:tests} is the test plan.
\end{abstract}

\tableofcontents

% ===========================================================================
\section{Why: the move is GPU-bound and does redundant work}
% ===========================================================================
Measured on zero\_polymer ($L=566$, zero\_polymer\_recap.tex Section~7):
\begin{center}
\begin{tabular}{@{}lr@{}}
\toprule
operation & time \\
\midrule
forward, no gradient                    & 60.3 ms \\
forward + backward (any number of relaxed sites) & 150.4 ms \\
current move: $2\times$(fwd+bwd) + 1 fwd & 361 ms (model), 364 ms (measured) \\
two jobs on the card, per job            & 789 ms \\
GPU utilisation, two jobs                & 100\% \\
\bottomrule
\end{tabular}
\end{center}
The card is saturated by one job, so running chains in parallel gains nothing
(two jobs give $\approx8\%$ less total throughput than one). The only way to go
faster is less GPU work per move.

The current move (\file{classes/rate\_sampler.py:\_step}) does three model
evaluations:
\begin{enumerate}
  \item gradient at $A$: forward+backward with the chosen sites softened,
        $\mathrm{softmax}(x/\Tsftm)$, every other site one-hot;
  \item gradient at $B$: the same, at the candidate, for the reverse proposal
        probability $a_{B\to A}$;
  \item exact energy at $B$: a forward pass without gradient on the one-hot
        sequence.
\end{enumerate}
Two facts make this redundant. Pass 3 is needed only because passes 1 and 2
run on softened inputs, so their energy is not the exact energy. And pass 1 at
the next move recomputes a gradient that pass 2 has just computed, at the same
sequence if $B$ was accepted or at $A$ itself if it was rejected; it cannot be
reused today only because it is restricted to the sites drawn in this move.

% ===========================================================================
\section{The gradient at the exact sequence}
% ===========================================================================
\subsection{Definition}
ESM3 takes the sequence as a probability matrix $p\in\mathbb{R}^{L\times K}$
($K=25$) and embeds it linearly,
\begin{equation}
  e = p\,W_{\mathrm{embed}}
\end{equation}
(\file{customs/custom\_esm/models/esm3.py:manual\_embed}, a
\texttt{tensordot}). Everything downstream is smooth, so $U(p)$ is a smooth
function on all of $\mathbb{R}^{L\times K}$, not only on the simplex and not
only away from its corners. Its gradient at the one-hot point $p=\sigma$ of the
current sequence,
\begin{equation}
  g_{ik} = \frac{\partial U}{\partial p_{ik}}\bigg|_{p=\sigma},
\end{equation}
is well defined and is computed by one forward pass on the exact one-hot
sequence with \texttt{requires\_grad} on $p$, followed by one backward pass. The
forward pass of that evaluation \emph{is} the exact energy $U(\sigma)$: same
input, same code path (\texttt{predict\_attention} under the same bfloat16
autocast) as \texttt{\_exact\_energy}. The continuous relaxation is still what
makes this possible: it is what lets ESM3 accept a real-valued $p$. What changes
is only where the derivative is evaluated: at the corner itself instead of at a
softened point next to it.

$g_{ik}$ is the first-order change in $U$ when a small weight is moved at site
$i$ toward residue $k$. Only differences $g_{ik}-g_{ij}$ within a site enter the
proposal (Section~\ref{sec:proposal}), so directions that leave the simplex do
not matter.

\subsection{Relation to the current gradient}
Today the gradient is taken with respect to the logits $x$ with
$p=\mathrm{softmax}(x/\Tsftm)$. By the chain rule, at one site,
\begin{equation}
  \frac{\partial U}{\partial x_k}
  = \frac{p_k}{\Tsftm}\Bigl(g_k - \textstyle\sum_m p_m g_m\Bigr).
\end{equation}
The logits are a mean-centred one-hot vector (\file{ExtendedProtein.expand}):
$1-1/K$ on the current residue and $-1/K$ on the other 24, a gap of exactly 1.
At $\Tsftm=0.1$ this gives
\begin{equation}
  p_{\mathrm{cur}} = \frac{1}{1+24e^{-10}} = 0.998911,\qquad
  p_{\mathrm{off}} = p_{\mathrm{cur}}\,e^{-10} = 4.535\times10^{-5}
\end{equation}
for every non-current residue. For any two competitors $k,j$ of the current
residue,
\begin{equation}\label{eq:scale}
  \frac{\partial U}{\partial x_k}-\frac{\partial U}{\partial x_j}
  = c\,(g_k-g_j),\qquad c=\frac{p_{\mathrm{off}}}{\Tsftm}=4.535\times10^{-4}.
\end{equation}
So the two gradients rank the competitors identically and differ by one
global factor $c$. (The current gradient is evaluated at the softened point,
where the chosen site carries 0.1\% weight off its residue, so equality is up to
that $O(10^{-3})$ perturbation of the input; the test of
Section~\ref{sec:tests} measures it.)

\subsection{All sites in one pass}
The backward pass gives $g_{ik}$ for every $i$ at once. Its cost does not depend
on how many sites carry a gradient: forward+backward takes 150.8--151.1\,ms for
1 to 566 relaxed rows, and memory differs by 0.2\,MB in 6.4\,GB
(\file{code/zero\_polymer/attention\_locality.txt}). Self-attention is global, so
the backward crosses all 48 layers at every position regardless. The all-site
gradient is therefore free, and it is what allows reuse: a gradient for every
site is valid whatever sites the next move draws.

% ===========================================================================
\section{The move}\label{sec:proposal}
% ===========================================================================
The chain carries a cache $(\sigma, U(\sigma), g(\sigma))$: the current sequence,
its exact energy and its all-site gradient. One move:
\begin{enumerate}
  \item Draw $n$ distinct sites uniformly (\texttt{n\_sites}; symmetric,
        cancels).
  \item At each chosen site $i$, with the cached $g=g(A)$:
        \begin{equation}
          \mu_{ik}=-\frac{\dt^2}{2\Mass}\,g_{ik},\qquad
          \sigma_{\mathrm{p}}=\dt\sqrt{T/\Mass},\qquad
          \delta x_{ik}=\mu_{ik}+\sigma_{\mathrm{p}}\,\xi_{ik},\ \ \xi_{ik}\sim\mathcal N(0,1),
        \end{equation}
        target $j_i=\arg\max_{k\in\mathcal C(i)}\delta x_{ik}$ over the 19
        canonical competitors, and $\log a_{A\to B}=\sum_i\log a_{i\to j_i}$ from
        \texttt{log\_pointing\_prob}, unchanged.
  \item Build $B$. \textbf{One forward pass} on $B$'s one-hot matrix with
        gradient enabled gives $U(B)$; \textbf{one backward pass} gives $g(B)$
        for all sites.
  \item $\log a_{B\to A}$ from $g(B)$ at the same sites, targets the old
        residues (unchanged code).
  \item Accept with probability $\min(1,e^{\ell})$,
        $\ell=-\bigl(U(B)-U(A)\bigr)/T+\log a_{B\to A}-\log a_{A\to B}$.
  \item If accepted, the cache becomes $(B,U(B),g(B))$. If rejected, it stays
        $(A,U(A),g(A))$, and $g(B)$ is discarded.
\end{enumerate}
Model work per move: one forward and one backward, $\approx150$\,ms, whatever
the outcome and whatever $n$. At chain start, and after a restart, the cache is
filled by one extra forward+backward at the initial sequence.

The rest of the machinery is unchanged: the pointing probability, the
competitor set, the $\sqrt{T}$ scaling of $\dt$, the \file{data.dat} columns,
checkpointing. Gradient storage is $L\times K$ floats, 57\,kB.

% ===========================================================================
\section{Exactness}
% ===========================================================================
\subsection{With a deterministic gradient}
$g(\sigma)$ is a function of the state alone, so $a_{A\to B}$ and $a_{B\to A}$
are well-defined proposal probabilities and $\ell$ is the standard
Metropolis--Hastings ratio. Detailed balance with respect to $\pi\propto e^{-U/T}$
follows exactly as for the current sampler. Caching changes when $g$ is
computed, not what it is.

\subsection{With a non-reproducible backward}
The backward pass is not bit-reproducible: two gradients at the same input
differ by 0.18\% and change the preferred substitution at 3 of 566 sites
(gradient\_method.tex Section~6.3). So the gradient is really a random variable
$g\sim q(\cdot\mid\sigma)$, and a cached $g(A)$ and a later recomputation at $A$
can differ. The kernel stays exact on the extended state $(\sigma,g)$ with
target $\pi(\sigma)\,q(g\mid\sigma)$. The forward transition
$(A,g_A)\to(B,g_B)$ has probability $a(A\to B;g_A)\,q(g_B\mid B)$ and the reverse
$a(B\to A;g_B)\,q(g_A\mid A)$, so
\begin{equation}
  \frac{\pi(B)\,q(g_B\mid B)\;a(B\to A;g_B)\,q(g_A\mid A)}
       {\pi(A)\,q(g_A\mid A)\;a(A\to B;g_A)\,q(g_B\mid B)}
  = \frac{\pi(B)\,a(B\to A;g_B)}{\pi(A)\,a(A\to B;g_A)} ,
\end{equation}
which is exactly the $\ell$ above: the $q$ factors cancel and never need to be
known. The $\sigma$-marginal of the extended target is $\pi$. The energy needs
no such argument: the forward pass is bit-reproducible (noise floor exactly 0,
\file{scan\_U\_am\_results.txt}), so the cached $U$ is the energy.

The current sampler, which recomputes $g(A)$ every move, is the same argument
plus an extra refresh of $g$ at fixed $\sigma$; the one-pass move simply skips
the refresh.

\subsection{Things that must stay true in the implementation}
\begin{itemize}
  \item The $g(A)$ used in step 2 must be the cached one, the gradient the
        chain actually holds, never a recomputation.
  \item $\lambda_S$ must be 0. The entropy term $-\sum p\log(p+\epsilon)$ has a
        gradient of order $\log\epsilon$ at a one-hot corner, which would
        swamp $g$. It is 0 in every run; the code should refuse
        $\lambda_S\ne0$ in one-pass mode.
  \item $\lambda_{\mathrm{structure\_ce}}>0$ works the same way: its logits come
        from the same forward on $p$, so $U$ and $g$ still come from one pass.
  \item On restart, recompute the cache from the restored sequence (valid by
        the extended-state argument), or save $g$ in the checkpoint for
        bit-identical replay.
\end{itemize}

% ===========================================================================
\section{Calibration of the step size}
% ===========================================================================
The informedness depends on $\dt$ only through
$\alpha_{kj}=\dt\,(\text{gradient gap})/(2\sqrt{\Mass T})$ (informedness.tex).
By Eq.~\eqref{eq:scale} the gap in the new gradient is $1/c$ times the old one,
so the same informedness needs
\begin{equation}
  \dt_{\mathrm{new}} = c\,\dt_{\mathrm{old}},\qquad c = 4.535\times10^{-4}.
\end{equation}
Predictions for zero\_polymer, from the current rule
$\dt_{\mathrm{old}}=63\sqrt{T/2}$:
\begin{center}
\begin{tabular}{@{}rrr@{}}
\toprule
$T$ & $\dt_{\mathrm{old}}$ & $\dt_{\mathrm{new}}$ predicted \\
\midrule
2    & 63    & 0.0286 \\
32   & 252   & 0.114 \\
107  & 460   & 0.209 \\
300  & 772   & 0.350 \\
\bottomrule
\end{tabular}
\end{center}
In the sweep script this is \texttt{DT\_AT\_T2}$\,=0.0286$ with \texttt{DT\_EXP}
unchanged. These are predictions. The $p_{\mathrm{match}}$ curve must be
re-measured with the new gradient, and in the bulk (MUTS$\approx51$) as well as
at the floor, since the floor measurement is known to mislead
(zero\_polymer\_recap.tex Section~4.3).

% ===========================================================================
\section{Expected gain, and what it does not change}
% ===========================================================================
At equal acceptance the speed-up is the cost ratio, $361/150\approx2.4\times$ per
move. Projected: $\approx0.15$\,s per move, $\approx2.4\times10^4$ moves per hour
per job, and the 3\,$T\times$4 chains$\times$20k-move study in $\approx10$ hours
instead of $\approx24$. Acceptance should not change, because the proposal
distribution is the same up to the $O(10^{-3})$ shift in the evaluation point
and the $\dt$ rescaling. That is the main thing the tests have to confirm.

What it does not change: the energy, the target distribution, the pointing
probability, how informed the proposal is, and the multi-site question (cost is
still flat in $n$).

% ===========================================================================
\section{Test plan}\label{sec:tests}
% ===========================================================================
Ordered cheapest first. Each test has a pass criterion fixed before it runs.
\begin{description}
  \item[T0 Gradient correctness.] At the reference and at a MUTS=51 sequence,
        compare $g_{ik}-g_{ij}$ with a central finite difference of $U$ along
        $p_{ik}-p_{ij}$ (fp32, no autocast), at 20 sites $\times$ 3 pairs.
        Pass: relative error $<1\%$.
  \item[T1 Same energy.] $U$ from the gradient-enabled forward equals
        \texttt{\_exact\_energy} bit for bit, at 50 random sequences.
  \item[T2 Same ranking, predicted scale.] At 30 sites, reference and MUTS=51:
        the new $g$ against today's \texttt{\_grad\_pass\_sites} over the 19
        competitors. Pass: cosine of the centred vectors $\ge0.999$, same
        argmax except at near-ties, fitted scale within 1\% of
        $c=4.535\times10^{-4}$.
  \item[T3 Exact kernel, toy energy.] Extend \file{test\_nsite\_kernel.py}:
        enumerate the one-pass kernel on the 3-letter systems for
        $n=1,2,3$ and check detailed balance and stationarity to machine
        precision; repeat with Gaussian noise added to every gradient
        evaluation (the non-reproducible case) and check stationarity of the
        $\sigma$-marginal by long runs. Also check that a run that restarts and
        recomputes the cache gives the same distribution.
  \item[T4 Exact Boltzmann, ESM3.] \file{validate\_boltzmann\_toy.py} in
        one-pass mode at $T=107$ with $\dt$ set for $p_{\mathrm{match}}\approx0.15$,
        16\,000 and 32\,000 samples. Pass: KL and TV in line with the current
        sampler's (0.0316 / 0.0787 at $n=16\,000$) and decaying as $1/n$ and
        $1/\sqrt n$.
  \item[T5 Informedness and $\dt$.] \file{test\_informedness\_vs\_dt.py} with the
        new gradient, at the floor and at MUTS=51; locate
        $p_{\mathrm{match}}=0.15$ and compare with the predicted $\dt=0.0286$ at
        $T=2$.
  \item[T6 Cost.] Time per move (wall-clock and CPU) and peak memory, one job
        alone and two jobs, GPU utilisation. Pass: $\le170$\,ms per move alone.
  \item[T7 Same dynamics.] Old and new sampler, same starts, at the study's
        three temperatures, 5000 moves each, several seeds: acceptance, $H_d$
        growth and $U$ per move must agree within seed-to-seed spread. Then
        equilibrium statistics (mean $U$, mean $H_d$, site entropy, $q$) after
        equilibration, by the n\_sites study's criterion.
  \item[T8 Speed where it counts.] Decorrelation per wall-second (the n\_sites
        study's score) for old and new. This is the thesis figure: the
        expected value is $\approx2.4\times$.
  \item[T9 Compatibility.] \texttt{n\_sites}$>1$, \texttt{--extend}, restart,
        \texttt{run\_equilibration.py}, and $\lambda_S\ne0$ refused.
\end{description}

\paragraph{Implementation sketch.} A pars key (e.g.\ \texttt{grad\_mode} =
\texttt{site\_soft} | \texttt{one\_pass}, default today's behaviour so every
old run replays), a method that returns $(U,g)$ from one pass on the one-hot
matrix, and the cache held in \texttt{sample()} across calls to
\texttt{\_step}.

\end{document}
